\chapter{Zusammenfassung}\label{chap:Zusammenfassung}
Diese Bachelorarbeit beschäftigt sich mit einem Ansatz, der die Gesamtlaufzeit einer Testsuite eines Softwareprojektes reduzieren soll, indem vermieden wird, dass Testfälle, die bereits für den aktuellen Stand der Softwareprojekts evaluiert wurden, erneut ausgeführt werden. Dieser Ansatz bewegt sich in einem Szenario, bei dem eine unbestimmte Teilmenge aller auszuführenden Testfälle bereits durchlaufen wurde. Beispielsweise kann es sich dabei um einen Build-Server handeln, der einen Commit von einem Entwickler erhält, welcher bereits bei sich lokal eine Teilmenge der Tests laufen lassen hat.

Um diesen Ansatz zu realisieren, wird ein Tool erstellt, welches möglichst unabhängig vom verwendeten Testframework entschieden kann, welche Testfälle bereits durchlaufen wurden, und welche nicht. Anhand dieser Evaluation wird dann nur eine Teilmenge der gesamten Testuite ausführt, wobei eine insgesamt 100\%-ige Abdeckung der Testsuite gewährleistet sein muss.\\
Dieses Tool soll sich in eine bestehende CI-Pipeline integrieren lassen, auch hier möglichst unabhägig von der verwendeten Software.\\
Hierfür werden im Folgenden die Anforderungen an ein Projekt, in das das Tool integriert werden kann, beschrieben und erklärt.\\ %, sowie unterschiedliche Testausgabeformate hinsichtlich ihrer Eignung miteinander verglichen und diskutiert.\\
